\documentclass[12pt]{article}
\usepackage[utf8]{inputenc}
\usepackage[T1]{fontenc}
\usepackage{hyperref}
\title{Phong-Kniang Language Resources for K-12 Curriculum Development}
\author{Compiled by Claude Sonnet 5.5 (claude-sonnet-5-5) for Andr\'as Kornai}
\date{}
\begin{document}
\maketitle

\section{Introduction}
Phong-Kniang (ISO 639-3 pnx; Glottocode phon1246) is an Austroasiatic (Khmuic) language of Houaphan Province, Laos, also called Phong, Kniang, Kha Phong or Khaniang. Published speaker estimates vary widely (about 30,700 on Wikipedia, about 1,000 on Verbix). It should not be confused with the Pong people who speak a Viet-Muong language. The search found no school materials in the language.

\section{Primary Educational Resources}
No primers, textbooks or school programmes in Phong-Kniang were found in the searches.

\section{Linguistic and Reference Materials}
\begin{itemize}
\item Kniang language (Wikipedia overview): \url{https://en.wikipedia.org/wiki/Kniang_language}
\item Phong Kniang numeral data (MPI-EVA Numeral Systems of the World): \url{https://lingweb.eva.mpg.de/channumerals/Phong-Kniang.htm}
\item Grambank entry for Phong-Kniang: \url{https://grambank.clld.org/languages/phon1246}
\item ASJP word-list entry: \url{https://asjp.clld.org/languages/PHONG_KNIANG}
\item Ksingmul people (related Khmuic community): \url{https://en.wikipedia.org/wiki/Ksingmul_people}
\end{itemize}

\section{Religious and Community Resources}
\begin{itemize}
\item ScriptureEarth page for pnx (the search result listed it; its content was not inspected): \url{https://www.scriptureearth.org/00eng.php?iso=pnx}
\item Joshua Project, Pong in Laos (note the distinction from the Viet-Muong Pong): \url{https://joshuaproject.net/index.php/people_groups/14442/LA}
\end{itemize}

\section{Digital and Audio Resources}
\begin{itemize}
\item Verbix endangered-languages page: \url{https://docs.verbix.com/EndangeredLanguages/PhongKniang}
\item Rosetta Panglossia entry: \url{https://rosettapanglossia.longnow.org/index.php?title=Pnx}
\end{itemize}

\section{Curriculum Development Context}
Little documentation exists beyond basic descriptions and word lists. There is no known orthography or schooling in the language; Lao is the likely school language. A curriculum would first need community-led orthography development and documentation.
\end{document}